\label{sec:related_work}

\textbf{Open LLM families.} Recent open model families provide the foundation
for local deployment studies. Qwen3 and Qwen2.5 report strong multilingual,
reasoning, and code capabilities across practical parameter scales
\cite{qwen3,qwen25}. DeepSeek LLM and DeepSeek-R1 emphasize open-source scaling
and reasoning-oriented post-training~\cite{deepseekllm,deepseekr1}, while Phi
models target compact local capability~\cite{phi3,phi4}. Mistral 7B remains an
important efficient 7B-class baseline~\cite{mistral7b}, and Gemma and StarCoder2
provide additional open-weight general and code-specialized references
\cite{gemma2024open,gemma2,gemma3,starcoder2}. These works show that
single-GPU deployment is plausible, but they do not by themselves answer which
model is best for a local railway-domain workload.

\textbf{LLM evaluation.} Benchmark suites such as MMLU, GSM8K, HumanEval, and
C-Eval measure broad knowledge, mathematical reasoning, code generation, and
Chinese-domain understanding~\cite{mmlu,gsm8k,humaneval,ceval}. Evaluation
frameworks such as the Language Model Evaluation Harness improve
reproducibility by standardizing tasks and metrics~\cite{lm_eval}. Broader
surveys caution that LLM evaluation remains sensitive to benchmark coverage,
prompting, leakage, and reporting practice~\cite{liang2022holistic,laskar2024systematic}. Our study follows this line but adds a fixed hardware
budget and local railway data, because leaderboard-style public scores are
insufficient for deployment decisions.

\textbf{Parameter-efficient adaptation.} Parameter-efficient transfer learning
and adapters reduce the cost of adapting pretrained models
\cite{houlsby2019parameter,mahabadi2021compacter,pfeiffer2020adapterfusion}.
LoRA injects trainable low-rank matrices while freezing base weights
\cite{lora}; QLoRA combines LoRA with 4-bit quantization and enables
fine-tuning larger models within much smaller memory budgets~\cite{qlora}.
Related work extends efficient adaptation to federated and heterogeneous
settings~\cite{zhang2023fedpetuning,yi2023fedlora,cho2023heterogeneous,cai2022fedadapter,yan2024federa,gao2024fedpt,chen2026federated}. This
literature motivates the QLoRA part of our workflow, but the present paper
evaluates QLoRA as a practical railway-domain adaptation tool rather than as a
new training algorithm.

\textbf{Compression, pruning, and efficient deployment.} Model compression and
pruning have a long history, from deep compression~\cite{han2015deep} to
movement pruning~\cite{sanh2020movement}, structural pruning~\cite{ma2023llm,xia2022structured,tao2023structured}, and one-shot sparse LLM pruning
\cite{frantar2023sparsegpt,sun2023simple}. Recent resource-efficient LLM
surveys organize these methods across compression, inference, serving, and
adaptation~\cite{bai2024beyond,yao2025deployment}. Empirical deployment studies
further show that edge and mobile environments require joint measurement of
memory, latency, throughput, and task quality~\cite{xiao2024mobile,dhar2024empirical,seymour2024small,wani2026benchmarking,zhao2024hetegen}.
This motivates our decision to report efficiency metrics alongside task scores.

\textbf{Resource-constrained applications.} Domain studies in education,
clinical information extraction, emergency medicine, low-resource translation,
and humanities research show that constrained environments are not only a
systems problem but also an application-quality problem~\cite{barros2025evaluating,builtjes2025leveraging,mensah2024ambulance,posada2024medical,li2025comparative,zhong2026opportunities}. These works support the premise that
local LLM studies should account for both deployment constraints and domain
task validity. Railway bilingual QA shares this pattern: the model must be
affordable to run locally and reliable enough for specialized terminology and
technical-document assistance.
